\section{怎样写个人计划和总结}

每当一个新学年开始或者一个新阶段开始，我们往往要订一个学习计划或者工作计划；每当一个学年或者一个阶段结束，我们也往往要作一个学习总结或者工作总结。这说明个人计划是预计在一定时间内要完成的任务、要做的工作的书面化、条理化和具体化，而个人总结则是我们对规定时限内的某项任务、某项工作的回顾、评价和结论。

古语说：凡事预则立，不预则废。这说明了计划的重要作用。而总结可以使我们全面、系统地了解学习或工作的情况，从中获得经验，吸取教训，所以有人说：总结就是提高，总结就是转化。由于写计划或写总结，都需要从实际出发，进行分析研究，由表及里地加以归纳，并且还要条理清楚地表达出来，这对训练我们思路清晰、善于分析归纳的逻辑思维能力和概括能力、表达能力是很有帮助的。

那么，我们怎样来写个人计划和总结呢？

首先谈谈怎样写个人计划。个人计划与文件计划不同，属于普通计划，可以自由、灵活地写，在行文上既可以成文叙述，又可以分条列项地表达。一般说来，它由标题、正文、结尾三个部分构成。

（1）标题：必须写明计划适用的时间和计划的性质，有时也可以把是谁的计划注明。例如：《本学期学习计划》，《我在“大战红五月”中的工作计划》。

（2）正文：包括前言和计划事项。

前言部分，要简明扼要地说明为什么制定这份计划，制订计划的依据是什么。文字要简明扼要。

计划事项是计划的主体，这一部分一般要分项来写，有时大的项目之下有小的项目（小的项目是在大的项目下要做的一件事情或一项项工作）。每一项都应写明做什么、怎么做，什么时候完成。

（3）结尾：要写明订计划的人和日期两项。

制定个人计划应该注意：第一是从实际出发，要根据个人实际情况，目标不能订得过高，以免计划落空；也不能过低，而使计划失去意义。以切实可行为好。第二是要订得具体明确。要使目的、任务、要求、方法、步骤、措施毫不含糊，便于执行。第三是既突出重点，又适当安排一般事项。一般说来，应把重点事项突出地放在首位，以免胡子眉毛一把抓；再把其他各项按其重要程度依次排列，以免“单打一”地只有重点而无一般。

我们再来谈谈怎样写个人总结。

我们写的个人总结，按性质分，有学习总结、思想总结、工作总结、生产总结、科研总结等；按时间分，有月份总结、季度总结、阶段总结、年终总结、学期总结等。总结的种类不同，内容各异，写法也多种多样。但不管有多少种写法，总结的正文一般都包括基本情况回顾、主要成绩（收获）、基本经验、存在问题（或失败教训）、今后意见（或努力方向）五个部分。总结的标题和末尾，跟计划的写法相同。

（1）基本情况回顾：这是总结的开头，它的作用相当于计划的前言部分。但是计划的前言有时可以省掉，总结的基本情况回顾一般不可省略。这部分，一般是概述基本情况，把这段时间里执行计划的经过情况、主要成绩作一个总括性的介绍，写法必须简明扼要，文字较少。

（2）主要成绩（收获）：这是总结的不可缺少的一个方面，往往概括为几点来写。只有看到成绩，才能鼓舞士气，以利再战。成绩要写得具体，要有事实材料和必要的数据加以印证，有时也可以做对比说明。文字不宜过多。

（3）基本经验：这是总结的重要部分，说明工作的做法、体会和取得成绩的原因。从分析研究客观事实和做法入手，找出规律性的东西；即是说，一方面要摆出具体的典型事例，另一方面要能从这些材料中提炼出观点、看法。既不能空泛议论，不写具体材料；又不能就事论事，忽略了理论分析。这样总结出来的经验，才是有指导作用的“真经”。在写法上，往往根据内容把经验分成几个方面，每个方面用一个论点作领句概括出一个方面的经验，下面陈述具体内容。这与议论文用论据来论证论点的正确性颇为相似，要做到材料与观点统一。

（4）存在问题：主要是写没有完成或没有做好的事情，或者有待于解决的问题，以及缺点、错误和产生这些缺点、错误的原因。找出存在的问题，引为教训，以利再战。这部分要写得明确，但文字不宜过多，概要写出就行了。

（5）今后意见：主要是写努力的方向和今后的打算。这部分也要写得明确，文字却比较少，用来结束全文。如果是专题性的经验总结，这一部分一般可以略去。

我们在写个人总结时，并不是一定要逐项写全。有的可以合并，例如主要成绩可以合并到基本情况中去写，经验和教训也可以合并到一块写；有的可以略去，例如存在的问题和今后的意见，在以总结经验为主的总结里，一般可以不写。这都要根据总结的性质和目的要求来定，由于侧重点的不同，总结的写法也可以灵活安排。

写总结应该注意：

第一，要突出主要问题。不少人写总结，总想面面俱到，事无巨细，样样皆写，生怕有所遗漏，结果弄得文字冗长，而且把主要之点淹没在泛泛而谈的汪洋大海里；还有些人在写总结时，只是罗列材料，公布“流水账”，看不出中心内容究竟是什么。这样的总结究竟有多大的指导意义呢？我们一定要根据实际情况和总结目的，确定中心内容，突出主要问题，观点鲜明地表达出来，做到在不长的篇幅中，使材料和观点有机地结合起来，取得应有的效果。

第二，要找出规律性的东西。写总结主要是为了找出以前工作中的基本经验，乃至应该吸取的教训，以指导今后的工作。如果只是看到一些现象，叙述一通过程，或者只根据一些偶然现象轻易地得出结论，却没有找出规律性的东西，就失去了用总结经验来指导今后工作的意义。所以，我们一定要在掌握必要材料的基础之上，认真地分析、研究，去粗存精，由表及里，把一些感性印象上升到理性认识上来，从中归纳出带有普遍性的道理，思考出带有规律性的经验。

第三，要真实，切合自己的实际情况。总结既是实践的回顾和概括，内容就该切切实实反映出自己的实际情况，而不能随意夸大或缩小，更不能无中生有地编造，必须实事求是；总结出的经验，要从实践当中来，既不能主观武断，也不能任意拔高，必须恰如其分；而这些实际情况或基本经验，应该带有“自己的”特色，是自己实践的总结，必须非同一般。如果自己写的个人总结，改头换面，说哪个人的都行，哪年都可用，取消了个人的特殊性，那么，可以肯定这是一篇失败的没有意义的总结。
\newpage